\chapter{接口与运行期多态}
\label{ch:13}

到第 12 章为止，你写的每个函数都只认一种具体类型：\texttt{int64} 也好，\texttt{Circle} 也好，参数的类型写死了，就只能吃那一种。

可现实里常有这种事：一个函数想同时处理「圆」和「正方形」，因为它俩都需要「算面积」。你不会想给每个类型都抄一份逻辑。这一章讲 Magel 怎么解决它——\texttt{interface}。

\section{问题：长得不一样，但都会做同一件事}

先看一个真实场景。你有一个圆、一个正方形，都要算面积：

\begin{lstlisting}[style=magel]
struct Circle { int64 r; };
struct Square { int64 side; };

int64 circleArea(Circle c) { return 3 * c.r * c.r; }
int64 squareArea(Square s) { return s.side * s.side; }
\end{lstlisting}

两个函数，逻辑几乎一样，只是字段名不同。再加一个三角形呢？再抄一份。

问题不在「代码重复」，而在没有共同的类型。\texttt{circle\allowbreak{}Area} 收不了 \texttt{Square}，\texttt{square\allowbreak{}Area} 收不了 \texttt{Circle}。你没法写「一个对所有图形都能算面积的函数」。

接口就是来解决这件事的：把「会算面积」这件事本身，变成一个类型。

\section{\texttt{interface}：给「能力」起个名字}

\begin{lstlisting}[style=magel]
interface Shape {
    int64 area();
    String describe();
};
\end{lstlisting}

这段在读作：

\begin{itemize}
\item \texttt{interface Shape} 声明了一个叫 \texttt{Shape} 的接口；
\item 里面两条方法声明（只有签名，没有函数体）：想要成为 \texttt{Shape}，就得有 \texttt{area(\allowbreak{})\allowbreak{}} 和 \texttt{describe(\allowbreak{})\allowbreak{}}；
\item 接口里只能有方法，不能有字段——这是硬规矩。
\end{itemize}

换句话说，\texttt{Shape} 描述的不是「有什么数据」，而是「会做什么」。这种类型，学名就叫契约（contract）：谁签了它，谁就得履行这两条。

\begin{tip}{提示}
\texttt{interface}、\texttt{struct}、\texttt{enum} 的声明在 \texttt{\}} 处结束，末尾的 \texttt{;} 可省略。接口中的抽象方法声明仍需要分号。
\end{tip}

\section{结构体实现接口}

写一个真的圆，让它「实现」\texttt{Shape}——写法是在结构体名字后面加 \texttt{:\allowbreak{} Shape}：

\begin{lstlisting}[style=magel]
interface Shape {
    int64 area();
    String describe();
};

struct Circle : Shape {
    int64 r;

    int64 area() { return 3 * r * r; }
    String describe() { return "圆"; }
};

int64 main() {
    Circle c;
    c.r = 2;
    println("{}，面积 {}", c.describe(), c.area());
    return 0;
}
\end{lstlisting}

\begin{lstlisting}[style=out]
圆，面积 12
\end{lstlisting}

\texttt{struct Circle :\allowbreak{} Shape} 的意思就是「\texttt{Circle} 承诺自己是 \texttt{Shape}」。它把 \texttt{Shape} 要的两条方法都写了，于是契约生效。

关键规矩：必须完全实现，一个都不能少。 少写一条，编译器直接拒绝：

\begin{lstlisting}[style=magel]
interface Shape {
    int64 area();
    String describe();
};

struct Bad : Shape {
    int64 r;
    int64 area() { return 3 * r * r; }
};

int64 main() { return 0; }
\end{lstlisting}

报错说得非常清楚：

\begin{lstlisting}[style=out]
magel：struct Bad 实现了接口 Shape，但没有完全实现：describe() -> std::lang::String
（接口要求同名、同返回类型、同形参类型；一个都不能少）
\end{lstlisting}

注意「同名、同返回类型、同形参类型」这三点——返回类型写错也算没实现。把 \texttt{area(\allowbreak{})\allowbreak{}} 写成返回 \texttt{String}，一样通不过。

\begin{tip}{提示}
一个细节：在 \texttt{main} 里 \texttt{c.\allowbreak{}describe(\allowbreak{})\allowbreak{}} 是静态调用——编译器早知道 \texttt{c} 是\texttt{Circle}，直接找到 \texttt{Circle:\allowbreak{}:\allowbreak{}describe} 就行，根本不需要接口。接口的价值要等到13.7 节，类型在运行期才确定的时候才显出来。
\end{tip}

\section{默认方法：接口可以带实现}

接口里的方法可以自带函数体，成为「默认实现」：

\begin{lstlisting}[style=magel]
interface Greeter {
    String hello() { return "你好，我是"; }
};

struct Dog : Greeter {
    String name() { return "旺财"; }
};

int64 main() {
    Dog d;
    println("{}{}", d.hello(), d.name());
    return 0;
}
\end{lstlisting}

\begin{lstlisting}[style=out]
你好，我是旺财
\end{lstlisting}

\texttt{Dog} 只写了 \texttt{name(\allowbreak{})\allowbreak{}}，\texttt{hello(\allowbreak{})\allowbreak{}} 直接用接口给的默认体。这在「大部分实现都一样、少数才特殊」时很省事——想让某个类改写默认行为，在结构体里重写一遍同名方法即可。

\begin{tip}{提示}
接口方法也可以有默认参数，但枚举载荷字段和接口方法声明这两处不支持默认参数（见语言参考第 6 节）。这一条初学阶段用不到，记个位置就行。
\end{tip}

\section{接口继承}

接口之间也能继承。\texttt{Named :\allowbreak{} Shape} 读作「\texttt{Named} 是更具体的 \texttt{Shape}」：

\begin{lstlisting}[style=magel]
interface Shape {
    int64 area();
};

interface Named : Shape {
    String name();
};

struct Square : Named {
    int64 side;

    int64 area() { return side * side; }
    String name() { return "正方形"; }
};

int64 main() {
    Square s;
    s.side = 4;
    println("{}，面积 {}", s.name(), s.area());
    return 0;
}
\end{lstlisting}

\begin{lstlisting}[style=out]
正方形，面积 16
\end{lstlisting}

\texttt{Square} 签名签的是 \texttt{Named}，但 \texttt{Named} 继承了 \texttt{Shape}，所以两个要求一起生效：\texttt{area(\allowbreak{})\allowbreak{}} 和 \texttt{name(\allowbreak{})\allowbreak{}} 都得有。

\section{一个结构体实现多个接口}

结构体可以同时签多个契约，用逗号隔开：\texttt{struct X :\allowbreak{} A,\allowbreak{} B}。如果两个接口要的是同一个签名的方法，只算一条要求——去重的：

\begin{lstlisting}[style=magel]
interface Shape {
    String describe();
};
interface Named {
    String describe();
    String name();
};

struct Circle : Shape, Named {
    int64 r;
    String describe() { return "圆"; }
    String name() { return "Circle"; }
};

int64 main() {
    Circle c;
    println("{} {}", c.describe(), c.name());
    return 0;
}
\end{lstlisting}

\begin{lstlisting}[style=out]
圆 Circle
\end{lstlisting}

\texttt{Shape} 和 \texttt{Named} 都要 \texttt{String describe(\allowbreak{})\allowbreak{}}，它们同名、同返回、同参数，所以 \texttt{Circle} 里只写一遍 \texttt{describe(\allowbreak{})\allowbreak{}} 就同时满足了两个接口。反过来，要是两个接口要的同名方法签名不一样，那就是冲突，写不出来。

\section{接口引用：多态真正登场的地方}

前面六节，接口看起来只是「一份要遵守的清单」——实现完就完了，调用还是静态的。这一节才是重点：接口名可以当引用的目标类型，而这时调用的具体实现，由运行期决定。

先看它长什么样：

\begin{lstlisting}[style=magel]
Shape mut& view;          // view 是「Shape 引用」
view.area();              // 调哪个 area()？运行期才知道
\end{lstlisting}

\texttt{Shape mut\&} 不是普通引用。它是个「胖引用」——里面装了两样东西：对象的地址和对象的类型号。所以 \texttt{view.\allowbreak{}area(\allowbreak{})\allowbreak{}} 可以按类型号去挑对应的实现。这就是运行期多态：同一行代码，喂不同的对象，跑不同的函数。

那怎么拿到这么个引用？只能从 \texttt{std:\allowbreak{}:\allowbreak{}lang:\allowbreak{}:\allowbreak{}Any} 里解包。为什么要这么绕？因为 \texttt{Shape mut\&} 需要类型号，而裸引用（\texttt{\&c} 那种）里没有类型号——语言不让你硬凑：

\begin{lstlisting}[style=magel]
Circle c;
Shape mut& v = &c;        // X 报错
\end{lstlisting}

报错会教你正确的路：

\begin{lstlisting}[style=out]
magel：接口引用 Shape mut& 只能从 obj 解包来，不能拿 Circle mut& 凑：
先 `o.set(具体对象);` 再 `Shape mut& q = o.get();`（接口引用是「地址 + 类型号」的胖引用）
\end{lstlisting}

\texttt{std:\allowbreak{}:\allowbreak{}lang:\allowbreak{}:\allowbreak{}Any} 是第 8 章讲的那个「类型擦除的盒子」——它能装下任何类型，并记住装进去的是谁。解包时按接口取，它就去校验：这个对象真的实现了 \texttt{Shape} 吗？

完整流程：存进盒子 \ensuremath{\rightarrow} 按接口取出来 \ensuremath{\rightarrow} 当接口用。

\begin{lstlisting}[style=magel]
interface Shape {
    int64 area();
    String describe();
};

struct Circle : Shape {
    int64 r;
    int64 area() { return 3 * r * r; }
    String describe() { return "圆"; }
};

struct Square : Shape {
    int64 side;
    int64 area() { return side * side; }
    String describe() { return "正方形"; }
};

String report(Shape mut& shape) {
    return shape.describe() + "：" + format("{}", shape.area());
}

int64 main() {
    Circle c; c.r = 2;
    Square s; s.side = 3;
    Any box1; Any box2;
    box1.set(&c);
    box2.set(&s);
    Shape mut& shape1 = box1.getReference();
    Shape mut& shape2 = box2.getReference();
    println("{}", report(shape1));
    println("{}", report(shape2));
    return 0;
}
\end{lstlisting}

\begin{lstlisting}[style=out]
圆：12
正方形：9
\end{lstlisting}

一步步拆：

\begin{enumerate}
\item \texttt{std:\allowbreak{}:\allowbreak{}lang:\allowbreak{}:\allowbreak{}Any box1;} —— 声明一个空盒子；
\item \texttt{box1.\allowbreak{}set(\allowbreak{}\&c)\allowbreak{};} —— 把 \texttt{c} 装进去。这里传的是 \texttt{\&c}（\texttt{c} 的引用），盒子不复制对象，只是记住「地址 + 它是 \texttt{Circle}」；
\item \texttt{Shape mut\& shape1 =\allowbreak{} box1.\allowbreak{}get\allowbreak{}Reference(\allowbreak{})\allowbreak{};} —— 按 \texttt{Shape} 把引用取出来。类型号校验就发生在这里：\texttt{Circle} 确实实现了 \texttt{Shape}，放行；要是取成没实现的接口，这里就报运行期错误；
\item \texttt{report(\allowbreak{}shape1)\allowbreak{}} —— \texttt{report} 只认 \texttt{Shape mut\&}，它压根不知道里面有圆还是正方形。\texttt{shape.\allowbreak{}area(\allowbreak{})\allowbreak{}} 到运行期才按类型号分派——这就是多态。
\end{enumerate}

同一个 \texttt{report}，一行代码，两段输出。 加第三种图形，\texttt{report} 一个字都不用改。

\section{该用谁：静态调用还是接口？}

刚学多态容易走极端，觉得「接口多高级，全用接口吧」。别。看这张表：

{\small\setlength{\tabcolsep}{4pt}
\begin{longtable}{>{\raggedright\arraybackslash}p{\dimexpr 0.459\linewidth-2\tabcolsep\relax}>{\raggedright\arraybackslash}p{\dimexpr 0.229\linewidth-2\tabcolsep\relax}>{\raggedright\arraybackslash}p{\dimexpr 0.292\linewidth-2\tabcolsep\relax}}
\toprule
场景 & 用什么 & 为什么 \\
\midrule
\endhead
你在编译期就知道具体类型 & 直接调用 & 快，清楚，别绕 \\
一个函数要吃多种类型，且类型在运行期才定 & 接口引用 & 只有它能做 \\
只是想让几个类型共享一段代码 & 先想想泛型（第 14 章） & 泛型在编译期展开，往往更合适 \\
\bottomrule
\end{longtable}
}

经验法则：只有「调用者不知道具体是谁」的时候，才需要接口引用。能静态调用就静态调用——这也是官方样例把静态和动态两条路摆在一起对比的原因。

\section{小结}

\begin{center}
\begin{tabular}{ll}
\toprule
你会了 & 关键点 \\
\midrule
声明接口 & \texttt{interface Shape \{ int64 area(\allowbreak{})\allowbreak{}; \};} 可有方法与关联类型声明，不能有值字段 \\
实现接口 & \texttt{struct Circle :\allowbreak{} Shape \{ … \}}，必须完全实现，签名一字不差 \\
默认方法 & 接口里的方法可以自带函数体 \\
接口继承 & \texttt{interface Named :\allowbreak{} Shape \{ … \}} \\
多接口 & \texttt{struct X :\allowbreak{} A,\allowbreak{} B}，同名同签名的要求去重 \\
接口引用 & \texttt{Shape mut\& view}，是「地址 + 类型号」的胖引用 \\
怎么拿 & \texttt{box.\allowbreak{}set(\allowbreak{}\&obj)\allowbreak{}; Shape mut\& v =\allowbreak{} box.\allowbreak{}get\allowbreak{}Reference(\allowbreak{})\allowbreak{};} \\
运行时多态 & 方法按对象类型号分派；调用者无需知道具体类型 \\
别滥用 & 能静态调用就别上接口 \\
\bottomrule
\end{tabular}
\end{center}

\section{下一步}

接口解决了「一个函数吃多种类型」，但它是运行期付账——胖引用、类型号、装箱解包，都有代价。

第 14 章的泛型解决同一个问题，却是编译期付账：代码为你用到的每种类型分别展开一份，跑起来没有运行期开销。学会这两样，你就能在自己的代码里做正确的取舍了。

翻到\hyperref[ch:14]{第 14 章 泛型}。
